\clearpage
\onecolumn

\appendices

\section{Configuração do Roteador ROT-ACME}

\begin{listing}[H]
\centering
\begin{minted}[
    bgcolor=bgcolor,
    frame=lines,
    breaklines=true,
    breakanywhere=true,
    fontsize=\footnotesize,
]{bash}

vyos@vyos:~$ configure

vyos@vyos:~# set interfaces ethernet eth0 address 172.24.254.2/30
vyos@vyos:~# set interfaces ethernet eth1 address 172.24.1.254/24 

vyos@vyos:~# set interfaces ethernet eth1 vif 20 description 'INTRANET'
vyos@vyos:~# set interfaces ethernet eth1 vif 20 address '172.24.20.254/24'
vyos@vyos:~# set interfaces ethernet eth1 vif 30 description 'DC'
vyos@vyos:~# set interfaces ethernet eth1 vif 30 address '172.24.30.254/24'

vyos@vyos:~# set protocols static route 0.0.0.0/0 next-hop 172.24.254.1

vyos@vyos:~# set service dhcp-server shared-network-name INTRANET subnet 172.24.20.0/24 default-router '172.24.20.254'
vyos@vyos:~# set service dhcp-server shared-network-name INTRANET subnet 172.24.20.0/24 dns-server '172.16.30.11'
vyos@vyos:~# set service dhcp-server shared-network-name INTRANET subnet 172.24.20.0/24 start '172.24.20.10' stop '172.24.20.100'

vyos@vyos:~# set service dhcp-server shared-network-name DC subnet 172.24.30.0/24 default-router '172.24.30.254'
vyos@vyos:~# set service dhcp-server shared-network-name DC subnet 172.24.30.0/24 dns-server '172.16.30.11'
vyos@vyos:~# set service dhcp-server shared-network-name DC subnet 172.24.30.0/24 start '172.24.30.10' stop '172.24.30.100'

vyos@vyos:~# set service snmp community grs authorization rw
vyos@vyos:~# set service snmp community grs network 172.24.0.0/16
vyos@vyos:~# set service snmp community grs client 172.24.1.254
vyos@vyos:~# set service snmp trap-target 172.24.1.254 community grs

vyos@vyos:~# commit
vyos@vyos:~# save
vyos@vyos:~# exit 

\end{minted}
\caption{Configuração de Roteador ROT-ACME}
\label{lst:rotacme}
\end{listing}


\section{Configuração do Roteador R1}

O Código \ref{lst:rotR1} apresenta a configuração realizada no roteador R1 utilizando o sistema operacional VyOS. A configuração contempla endereçamento IP, NAT, protocolo OSPF e monitoramento via SNMP.

\begin{listing}[H]
\centering
\begin{minted}[
    bgcolor=bgcolor,
    frame=lines,
    breaklines=true,
    breakanywhere=true,
    fontsize=\footnotesize,
]{bash}

vyos@vyos:~$ configure
vyos@vyos:~# set interfaces ethernet eth0 address dhcp
vyos@vyos:~# set interfaces ethernet eth1 address 192.168.15.17/28
vyos@vyos:~# set interfaces ethernet eth2 address 192.168.15.33/28
vyos@vyos:~# set interfaces loopback lo address 10.1.1.1/32

vyos@vyos:~# set nat source rule 100 outbound-interface 'eth0'
vyos@vyos:~# set nat source rule 100 source address '192.168.15.0/24'
vyos@vyos:~# set nat source rule 100 translation address 'masquerade'

vyos@vyos:~# set protocols ospf area 0 network 192.168.122.0/24
vyos@vyos:~# set protocols ospf area 0 network 192.168.15.16/28
vyos@vyos:~# set protocols ospf area 0 network 192.168.15.32/28

vyos@vyos:~# set protocols ospf default-information originate 
vyos@vyos:~# set protocols ospf default-information originate metric 10
vyos@vyos:~# set protocols ospf default-information originate metric-type 2

vyos@vyos:~# set protocols ospf log-adjacency-changes
vyos@vyos:~# set protocols ospf parameters router-id 10.1.1.1
vyos@vyos:~# set protocols ospf redistribute connected metric-type 2
vyos@vyos:~# set protocols ospf redistribute connected route-map CONNECT

vyos@vyos:~# set policy route-map CONNECT rule 10 action permit
vyos@vyos:~# set policy route-map CONNECT rule 10 match interface lo

vyos@vyos:~# set service snmp community grs authorization rw
vyos@vyos:~# set service snmp community grs network 192.168.15.0/24
vyos@vyos:~# set service snmp community grs client 192.168.15.33
vyos@vyos:~# set service snmp trap-target 192.168.15.33 community grs

vyos@vyos:~# commit
vyos@vyos:~# save
vyos@vyos:~# exit

\end{minted}
\caption{Configuração do roteador R1 em VyOS}
\label{lst:rotR1}
\end{listing}

\section{Configuração do Roteador R2}

\begin{listing}[H]
\centering
\begin{minted}[
    bgcolor=bgcolor,
    frame=lines,
    breaklines=true,
    breakanywhere=true,
    fontsize=\footnotesize,
]{bash}

vyos@vyos:~$ configure

vyos@vyos:~# set interfaces ethernet eth0 address dhcp
vyos@vyos:~# set interfaces ethernet eth1 address 192.168.15.18/28
vyos@vyos:~# set interfaces ethernet eth2 address 192.168.15.35/28
vyos@vyos:~# set interfaces ethernet eth3 address 192.168.15.50/28
vyos@vyos:~# set interfaces loopback lo address 10.2.2.2/32

vyos@vyos:~# set nat source rule 100 outbound-interface name 'eth0'
vyos@vyos:~# set nat source rule 100 source address '192.168.15.0/24'
vyos@vyos:~# set nat source rule 100 translation address 'masquerade'

vyos@vyos:~# set protocols ospf area 0 network 192.168.15.0/28
vyos@vyos:~# set protocols ospf area 0 network 192.168.15.16/28
vyos@vyos:~# set protocols ospf area 0 network 192.168.15.32/28
vyos@vyos:~# set protocols ospf area 0 network 192.168.15.48/28

vyos@vyos:~# set protocols ospf default-information originate
vyos@vyos:~# set protocols ospf default-information originate metric 100
vyos@vyos:~# set protocols ospf default-information originate metric-type 2

vyos@vyos:~# set protocols ospf log-adjacency-changes
vyos@vyos:~# set protocols ospf parameters router-id 10.2.2.2
vyos@vyos:~# set protocols ospf redistribute connected metric-type 2
vyos@vyos:~# set protocols ospf redistribute connected route-map CONNECT

vyos@vyos:~# set policy route-map CONNECT rule 10 action permit
vyos@vyos:~# set policy route-map CONNECT rule 10 match interface lo

vyos@vyos:~# set service snmp community grs authorization rw
vyos@vyos:~# set service snmp community grs network 192.168.15.0/24
vyos@vyos:~# set service snmp community grs client 192.168.15.50
vyos@vyos:~# set service snmp trap-target 192.168.15.50 community grs

vyos@vyos:~# commit
vyos@vyos:~# save
vyos@vyos:~# exit 

\end{minted}
\caption{Configuração de Roteador R2}
\label{lst:rotR2}
\end{listing}

\section{Configuração do Roteador R3}

\begin{listing}[H]
\centering
\begin{minted}[
    bgcolor=bgcolor,
    frame=lines,
    breaklines=true,
    breakanywhere=true,
    fontsize=\footnotesize,
]{bash}
configure

# Interfaces
set interfaces ethernet eth0 address 192.168.15.1/28
set interfaces ethernet eth1 address '172.24.40.254/24'
set interfaces ethernet eth2 address 192.168.15.34/28
set interfaces ethernet eth3 address 192.168.15.49/28
set interfaces loopback lo address 10.3.3.3/32

# OSPF
set protocols ospf area 0 network 192.168.15.0/28
set protocols ospf area 0 network '172.24.40.0/24'
set protocols ospf area 0 network 192.168.15.32/28
set protocols ospf area 0 network 192.168.15.48/28

vyos@vyos:~# set protocols ospf log-adjacency-changes
vyos@vyos:~# set protocols ospf parameters router-id 10.3.3.3
vyos@vyos:~# set protocols ospf redistribute connected metric-type 2
vyos@vyos:~# set protocols ospf redistribute connected route-map CONNECT

set service dhcp-server shared-network-name TESTE subnet 172.24.40.0/24 default-router '172.24.40.254'
set service dhcp-server shared-network-name TESTE subnet 172.24.40.0/24 dns-server '172.16.4.2'

set service dhcp-server shared-network-name TESTE subnet 172.24.40.0/24 range 0 start '172.24.40.10'
set service dhcp-server shared-network-name TESTE subnet 172.24.40.0/24 range 0 stop '172.24.40.30'

# Route-map
set policy route-map CONNECT rule 10 action 'permit'
set policy route-map CONNECT rule 10 match interface 'lo'

# SNMP
set service snmp community grs authorization 'rw'
set service snmp community grs network '192.168.15.0/24'
set service snmp community grs client '192.168.15.1'
set service snmp trap-target '192.168.15.1' community 'grs'

commit
save
exit

\end{minted}
\caption{Configuração de Roteador R3}
\label{lst:rotR3}
\end{listing}

\section{Configuração de Switch EXOS1}

\begin{listing}[H]
\centering
\begin{minted}[
    bgcolor=bgcolor,
    frame=lines,
    breaklines=true,
    breakanywhere=true,
    fontsize=\footnotesize,
]{bash}

create vlan INTRANET tag 20
create vlan DC tag 30

configure vlan INTRANET add ports 2 untagged
configure vlan INTRANET add ports 3 untagged
configure vlan INTRANET add ports 4 untagged

configure vlan DC add ports 5 untagged

configure vlan INTRANET add ports 1 tagged
configure vlan DC add ports 1 tagged


configure snmp add community readwrite grs
configure snmp sysContact EXOS1
configure snmp sysName EXOS1
enable snmp access snmp-v1v2c 

save 

\end{minted}
\caption{Configuração de Switch EXOS1}
\label{lst:exos1}
\end{listing}

\section{Configuração de Switch EXOS2}

\begin{listing}[H]
\centering
\begin{minted}[
    bgcolor=bgcolor,
    frame=lines,
    breaklines=true,
    breakanywhere=true,
    fontsize=\footnotesize,
]{bash}

create vlan INTRANET tag 20
create vlan DC tag 30

configure vlan DC add ports 2 untagged
configure vlan DC add ports 3 untagged
configure vlan DC add ports 4 untagged
configure vlan DC add ports 5 untagged
configure vlan DC add ports 6 untagged

configure vlan INTRANET add ports 7 untagged

configure vlan INTRANET add ports 1 tagged
configure vlan DC add ports 1 tagged

configure snmp add community readwrite grs
configure snmp sysContact EXOS2
configure snmp sysName EXOS2
enable snmp access snmp-v1v2c 

save 



\end{minted}
\caption{Configuração de Switch EXOS2}
\label{lst:exos2}
\end{listing}

\section{Configuração de Switch DMZ}

\begin{listing}[H]
\centering
\begin{minted}[
    bgcolor=bgcolor,
    frame=lines,
    breaklines=true,
    breakanywhere=true,
    fontsize=\footnotesize,
]{bash}

configure vlan default del ports all

create vlan DMZ
configure vlan DMZ tag 10

configure vlan 10 add port 2,3,4 untagged
configure vlan 10 add port 1 tagged

configure vlan DMZ ipaddress 172.24.10.2/24
configure iproute add default 172.24.10.1

configure snmp add community readwrite grs
configure snmp sysContact EXOSDMZ
configure snmp sysName EXOSDMZ
enable snmp access snmp-v1v2c 

save 

\end{minted}
\caption{Configuração de Switch EXOS3}
\label{lst:exos3}
\end{listing}

\section{Arquivo monitor.py}

\begin{listing}[H]
\centering
\begin{minted}[
    bgcolor=bgcolor,
    frame=lines,
    breaklines=true,
    breakanywhere=true,
    fontsize=\footnotesize,
]{python}

ultima_not = {}
cooldown = 30

with open("/var/log/snort/alert.fast", "r") as f:
    f.seek(0, 2)

    while True:
        agora = time.time()

        linha = f.readline()

        if not linha:
            time.sleep(0.5)
            continue

        match = re.search(
            r"(\d+\.\d+\.\d+\.\d+).*->\s*(\d+\.\d+\.\d+\.\d+)",
            linha
        )

        if not match:
            continue

        ip_origem = match.group(1)
        ip_destino = match.group(2)

        if ip_destino != myIp:
            continue

        tipo = None

        if "ICMP DETECTADO" in linha:
            tipo = "ICMP"

        elif "SSH DETECTADO" in linha:
            tipo = "SSH"

        elif "HTTP DETECTADO" in linha:
            tipo = "HTTP"

        elif "HTTPS DETECTADO" in linha:
            tipo = "HTTPS"

        if not tipo:
            continue

        chave = f"{ip_origem}_{tipo}"

        if (
            chave not in ultima_not
            or agora - ultima_not[chave] > cooldown
        ):

            subprocess.run([
                "notify-send",
                "Snort",
                f"Ataque {tipo} detectado de {ip_origem}"
            ])

            ultima_not[chave] = agora
			

\end{minted}
\caption{Arquivo monitor.py}
\label{lst:monitor}
\end{listing}

